% The spoken body of the homily for the Twenty-fourth and Last Sunday after
% Pentecost, said on Sunday 22 November 2026. Continuous preaching, ready to
% read aloud after the Gospel: no heading, no direction to the preacher, no
% citation inside the speech. The loci, the relation to the reviewed
% interpretations, the audience, the word count and the pace estimate stand in
% the terminal note, which no one reads out. The blank lines are the paragraph
% breaks a speaker recovers his place by. The speech tells the scene on the
% Mount of Olivet and the disciples' question; gives the commands of flight as
% St John Chrysostom reads them, of the Roman war, and St Bede's account of the
% flight to Pella; draws Chrysostom's conclusion from what is fulfilled to
% what is to come; takes the false christs and the lightning with Chrysostom;
% meets the difficulty of "this generation" with Chrysostom; gives St Gregory
% on heaven and earth and the word that does not pass; hears the Gradual's past
% tense with St Augustine and the Epistle's thanks with Chrysostom; joins the
% Introit's promise to prayer with St Jerome; reads the Communion with the
% verse after it in Mark and Bede's condition; gives one response from that
% verse; and closes on the Secret and Gregory's gladness at the world's end.
% The honorific is tied to the name so that "St" never ends a line.
\raggedcolumns

Jesus had come out of the temple, and his disciples came to show him its
buildings. He told them that there would not be left there a stone upon a
stone. Afterwards, as he sat on the Mount of Olivet, they came to him
privately and asked him two things: when these things would be, and what would
be the sign of his coming and of the end of the world. Today's Gospel is taken
from his answer. It stops just before he says, ``But of that day and hour no
one knoweth.''

The first part of his answer is about Judea. When they see ``the abomination
of desolation, which was spoken of by Daniel the prophet, standing in the holy
place,'' those in Judea are to flee to the mountains. An abomination is a
thing God abhors, and this one, standing in the temple, leaves it desolate.
The man on the housetop is not to come down for anything in his house, and the
man in the field is not to go back for his coat. They are to pray that their
flight be not in the winter or on the sabbath. St~John Chrysostom took these
commands as a warning to the Jews of that time, before the Roman war, and he
found the proof in the sabbath. The apostles were not bound to keep the
sabbath, and they would not be in Jerusalem when the Roman general Vespasian
came.

The warning was obeyed. St~Bede records that when the Roman war drew near, the
Christians of that province were warned to leave, and they crossed the Jordan
to the city of Pella and stayed there for a time. For us, who hear this Gospel
long afterwards, its first part is a prophecy already fulfilled. What remains
is the coming of the Son of man on the clouds of heaven, when he will send his
angels to gather his elect from the four winds. Chrysostom tells anyone who
doubts it to test the Lord's sayings, and when he has found them true, ``from
what is past, let him believe also the things to come.''

The Lord gives no date, and he warns against false christs and false
prophets. ``Then if any man shall say to you, Lo here is Christ, or
there: do not believe him.'' And then: ``For as lightning cometh out of the
east and appeareth even into the west: so shall also the coming of the Son of
man be.'' Chrysostom explains the comparison. Lightning needs no herald,
because it shows itself everywhere at once.

Then comes the parable of the fig tree, and two sayings. The first is, ``Amen
I say to you that this generation shall not pass till all these things be
done.'' This saying has a difficulty. If ``this generation'' means the people
who sat with him on the Mount of Olivet, they all died, and none of them saw
the Son of man coming on the clouds. Chrysostom does not read it so. The Lord,
he says, is speaking ``not of the generation then living, but of that of the
believers,'' as a psalm speaks of ``the generation of them that seek him.''
Over the generation of the faithful, he says, ``shall nothing prevail, not
famine, not pestilence.'' Everyone who believes belongs to that generation,
and the promise that it will not pass is made to you.

The second saying is, ``Heaven and earth shall pass: but my words shall not
pass.'' St~Gregory the Great saw something strange in it. Nothing in the
bodily world lasts longer than heaven and earth, and nothing passes more
quickly than speech. A word is not complete until it has been spoken, and once
it has been spoken it is gone. The Lord's meaning, Gregory says, is that
whatever seems lasting to us does not last for ever without change, while his
speech, which passes, utters sentences that remain without change.

After the Epistle the Gradual gave thanks: ``Thou hast delivered us, O Lord,
from them that afflict us.'' In the psalm those words come from, the people go
on to complain that God has cast them off. St~Augustine, preaching on this
verse, says that it speaks of the future in the form of the past, because what
it promises is as certain as if it had been done. To us, he says, what is past
is certain and what is to come is uncertain. But he asks us to think of a
prophet to whom the future is as certain as the past is to us. What we
remember has happened, and nothing can now undo it; what the prophet knows
will be cannot fail to be.

St~Paul's prayer in the Epistle has both tenses. He asks for what the
Colossians do not yet have, that they may walk worthy of God and be
strengthened ``in all patience and longsuffering with joy.'' And he gives
thanks to the Father, ``who hath delivered us from the power of darkness and
hath translated us into the kingdom of the Son of his love,'' in whom we have
the remission of sins. Chrysostom notices that Paul, about to exhort them,
speaks of what is already theirs and not of what will be laid up for them.
These, he says, are the causes of what is to come. The forgiveness of your
sins and your place in the kingdom of the Son have been given already, and
they are the ground of everything you are still waiting for.

The Introit at the beginning of this Mass took its words from a letter that
the prophet Jeremias sent to the exiles in Babylon. God says to them, ``I
think thoughts of peace, and not of affliction; you shall call upon Me and I
will hear you; and I will bring back your captivity from all places.'' The
same letter tells them that the return will come when seventy years are
accomplished. The promise was certain, and still they were told to call on
him. St~Jerome, commenting on the letter, says that the Lord could certainly
have done what he had promised without the captives' prayer. Yet the Lord
urges them to pray, he says, so that they may be worthy to receive it. So the
certainty of God's word is no reason for us to stop asking for what he has
promised.

The Communion antiphon puts the promise and the prayer in one sentence. It is
the Lord's own word, from St~Mark: ``Amen, I say to you, whatsoever you ask
when you pray, believe that you shall receive, and it shall be done unto
you.'' Taken alone, it sounds as though it promised anything whatever to
anyone who believes. In St~Mark the Lord goes straight on: ``And when you
shall stand to pray, forgive, if you have aught against any man: that your
Father also, who is in heaven, may forgive you your sins.'' St~Bede, on that
verse, says that those who are not yet perfect should ask for the forgiveness
of their sins. Without doubt, he says, they will obtain what they ask,
provided that they first forgive others who sin against them.

Before Communion today, then, forgive anyone you hold something against, and
ask God to forgive your sins, believing that you shall receive what you ask.

At the offering the priest will ask God to turn all our hearts to him, ``that,
being delivered from earthly desires, we may pass on to the enjoyments of
heaven.'' St~Gregory says that those who love God are bidden to rejoice and be
glad at the world's end, because they soon find him whom they love, while what
they did not love passes away. You need not hear this Gospel of the end with
dread. You do not know the day, but you know that he will come, because his
words do not pass. If your heart is turned to him, that day will be the day
you find him whom you love.
